% from ega4-9-fr.tex
\oldpage[IV-3]{58}
  sauf la fibre générique $X_\eta$ sont projectives sur $k(s)$ (tous les
  préschémas $X_s$ étant de dimension 2).
\end{enumerate}
\end{remarks}

\begin{proposition}[9.2.3]
\label{IV.9.2.3-fr}
Soient $\mathbf P$ une propriété constructible (resp. ind-constructible) de
préschémas algébriques, $S$ un préschéma, $X$ un préschéma de présentation
finie sur $S$, $\mathcal F$ un $\mathcal O_X$-Module quasi-cohérent de
présentation finie. Alors l'ensemble $E$ des $s\in S$ tels que
$\mathbf P(X_s,\mathcal F_s,k(s))$ soit vraie est localement constructible
(resp. ind-constructible). En outre, si $S$ est irréductible de point générique
$\eta$, un des deux ensembles $E$, $S-E$ est un voisinage de $\eta$ dans $S$
(resp. $E$ est un voisinage de $\eta$ s'il contient ce point).
\end{proposition}

Pour démontrer ces assertions, on peut se borner au cas où
$S=\operatorname{Spec}(A)$ est affine. On sait alors qu'il existe un
sous-anneau $A_0$ de $A$ qui est une $\mathbf Z$-algèbre de type fini, un
$A_0$-préschéma de type fini $X_0$ et un $\mathcal O_{X_0}$-Module cohérent
$\mathcal F_0$ tels que $X$ soit isomorphe à $X_0\otimes_{A_0}A$ et
$\mathcal F$ à $\mathcal F_0\otimes_{A_0}A$ (8.9.1). Soit
$p:S\to S_0=\operatorname{Spec}(A_0)$ le morphisme correspondant à
l'injection $A_0\to A$, et soit $E_0$ l'ensemble des $s_0\in S_0$ tels que
\[
\mathbf P((X_0)_{s_0},(\mathcal F_0)_{s_0},k(s_0))
\]
soit vraie~; alors, en vertu de (9.2.2, (iv)), on a $E=p^{-1}(E_0)$~; on peut
par suite (1.8.2) se borner au cas où $S$ est le spectre d'une
$\mathbf Z$-algèbre de type fini, donc un schéma noethérien. Utilisons le
critère de constructibilité $(\mathbf 0_{\mathrm{III}},9.2.3)$ (resp. le
critère d'ind-constructibilité (1.9.10))~; on est alors ramené, en utilisant
comme ci-dessus (9.2.2, (iv)) et en remplaçant $S$ par un sous-schéma fermé
intègre de $S$, au cas où $S$ est noethérien et intègre, et où il faut prouver
que $E$ est rare dans $S$ ou contient un ouvert non vide de $S$ (resp. que $E$
contient un ouvert non vide de $S$ s'il contient le point générique)~; mais
cela est garanti en vertu de la condition (9.2.1, (ii)).

On notera que l'on n'a utilisé (9.2.1, (ii)) que lorsque $S$ est le spectre
d'un anneau intègre de type fini sur $\mathbf Z$. Il est clair d'autre part
que l'énoncé de (9.2.3) s'applique aussi lorsqu'il figure dans $\mathbf P$
plusieurs préschémas, Modules sur ces préschémas, morphismes de préschémas ou
homomorphismes de Modules. Il s'applique encore lorsqu'il y figure des parties
(en nombre fini) des préschémas considérés, pourvu qu'on impose à ces parties
la condition d'être \emph{localement constructibles}. En effet, la restriction
au cas où $S$ est affine montre qu'on peut se borner au cas où ces parties sont
\emph{constructibles}~: on applique alors (8.3.11) qui montre (avec les
notations précédentes) qu'une partie constructible de $X$ est l'image
réciproque d'une partie constructible de $X_0$ pour un choix convenable de
$A_0$.

\begin{corollary}[9.2.4]
\label{IV.9.2.4-fr}
Soient $\mathbf P$ une propriété constructible (resp. ind-constructible) de
préschémas algébriques, $X$, $Y$ deux $S$-préschémas de présentation finie,
$f:X\to Y$ un $S$-morphisme. Pour tout $s\in S$, posons
$X_s=X\times_S\operatorname{Spec}(k(s))$,
$Y_s=Y\times_S\operatorname{Spec}(k(s))$,
$f_s=f\times1:X_s\to Y_s$. Alors l'ensemble $E$ des $s\in S$ tels que pour
tout $y\in Y_s$, la propriété $\mathbf P(f_s^{-1}(y),k(y))$ soit vraie, est
localement constructible (resp. ind-constructible).
\end{corollary}

En effet, soit $Z$ l'ensemble des $y\in Y$ tels que
$\mathbf P(f^{-1}(y),k(y))$ soit vraie. Comme les fibres $f^{-1}(y)$ et
$f_s^{-1}(y)$ sont isomorphes, on voit que $E$ est l'ensemble des $s\in S$
tels que $Y_s\subset Z$~; si $g:Y\to S$ est le morphisme structural, on a
donc $E=S-g(Y-Z)$.

